\subsection{The global attainable law}
\begin{theorem}[Collision-uniform checkpoint classification]\label{thm:intrinsic-checkpoint}
Under the assumptions of this section there are $0<c\le C<\infty$
such that, for every $a\in K$ and integer $M\ge1$,
\begin{equation}\label{eq:intrinsic-checkpoint}
 c\Psi_{n,m}(M,a)\le\overline R_{M;n,m}^{a}
 \le R_{M;n,m}^{a,\max}\le C\Psi_{n,m}(M,a).
\end{equation}
The unconditional lower law uses independent uniform acquisition commands.
The upper code is deterministic, has at most $M$ reachable representatives,
and works on every admitted history. The lower law permits the independent
coding randomization of Definition~\ref{def:finite-state}. Constants may
depend on $K,C,\eta,\mu,n,m$, but not on $a$ or $M$.
\end{theorem}
\begin{proof}
Let $D_a=\operatorname{diag}(d_j)$. For one report word each coordinate
of $z_a(h)$ has the form $N_j(u)/Z(u)$, with numerator and evidence
polynomials of degree at most $n$ in the command variables and
$Z(u)\ge\kappa^n$. Thus
\[
 S_a=\{D_az_a(h):|h|=n\}
\]
is compact and has coefficient-independent bounded semialgebraic format:
it is a finite union of projections of the formulas
$Z(u)y_j=d_jN_j(u)$ on the command cube. Prior integrals and node
values are real coefficients in these formulas. They need not depend
semialgebraically on $a$.

Individually normalizing the $n$ report factors puts them in affine
hyperplanes of dimension $r-1$. The moment map is rational in those
coordinates. Hence $\dim S_a\le p$. Lemma~\ref{lem:newton-attainment}
also gives $|z_{a,j}|\le B$, so $S_a$ lies in the rectangle with
half-widths $Bd_1\ge\cdots\ge Bd_{q_m}$. By
Lemma~\ref{lem:tame-rectangle}, it has at most $M$ reachable
representatives with covering radius bounded by
\[
 C\max_{\ell\le p}
       \left(M^{-1}\prod_{j\le\ell}d_j\right)^{1/\ell}.
\]
The physical metric equivalence follows from \eqref{eq:leja-metric}
and the full-column-rank product-probe matrix $G_m$. A fixed tie rule
makes nearest-representative encoding Borel. Squaring and applying
\eqref{eq:leja-volume} gives the upper bound in
\eqref{eq:intrinsic-checkpoint}, with the integer budget already enforced
by Lemma~\ref{lem:tame-rectangle}.

For the lower bound, first average any randomized decoder conditional on
its label; squared loss cannot increase. Orthogonally project its query
centers onto the affine physical query span, apply a fixed left inverse
of $G_m$, and then the inverse of the leading nonzero block of $L_a$.
These maps have uniformly bounded norms. Projection onto any first
$\ell\le\min(p,s)$ coordinates reduces the squared-error problem to
the marginal of the cube in Lemma~\ref{lem:newton-attainment}, scaled
by $d_1,\ldots,d_\ell$. The probability that $M$ radius-$b$ balls
cover this uniform rectangle is at most
\[
 C_\ell M b^\ell/(d_1\cdots d_\ell).
\]
Choose $b=c_\ell(d_1\cdots d_\ell/M)^{1/\ell}$ with the
constant small enough that at least half the minorized mass remains.
Its mean squared error is at least a fixed multiple of $b^2$.
An encoder's randomized assignment cannot improve upon the distance to
the nearest of its centers. Independent public coding randomness can be
fixed and then averaged, since it does not generate the acquisition law.
Maximize over $\ell$. For $\ell>s$ the determinant volume is zero,
so no additional term is needed. Equation~\eqref{eq:leja-volume}
now gives the lower bound. The unconditional optimum is bounded above
by the minimax optimum, completing the proof.
\end{proof}

\begin{theorem}[Causal classification across all collision strata]\label{thm:intrinsic-streaming}
For the same fixed experiment and horizon, there exist $c,C>0$ such that
for every known $a\in K$ and every integer $M\ge1$ there exists a
\emph{single, stage-compatible} deterministic $M$-state transducer with
maximum-history, maximum-checkpoint regret at most $C\Xi_N(M,a)$.
Every $M$-state transducer has maximum unconditional checkpoint regret at
least $c\Xi_N(M,a)$ under one common independent uniform exploration law.
Equivalently,
\begin{equation}\label{eq:intrinsic-streaming}
 c\Xi_N(M,a)\le\mathfrak R_{M,N}^{a,\mathrm{av}}
 \le\mathfrak R_{M,N}^{a,\max}\le C\Xi_N(M,a).
\end{equation}
The transducer may depend on $a$ and $M$. Constants are uniform in both;
codebooks are not asserted to be calibration-blind. It stores only its
representative index. The clock, known calibration and fixed real-valued
program are read-only, as in Definition~\ref{def:finite-state}.
\end{theorem}
\begin{proof}
At stage $n$ let $m=N-n$ and store a representative in the reachable
\emph{raw} moment image indexed by $\Lambda_m$, including the constant
coordinate $v_{me_0}=1$. For a new factor
$f(t)=\sum_i f_i(g,x)t^{a_i}$, the exact update is
\begin{equation}\label{eq:formal-multiindex-update}
 v'_\alpha=
 \frac{\sum_i f_i(g,x)v_{\alpha+e_i}}
      {\sum_i f_i(g,x)v_{e_i+(m-1)e_0}},
 \qquad \alpha\in\Lambda_{m-1}.
\end{equation}
Every index on the right belongs to $\Lambda_m$. Coincident labels
have equal moment values and cause no ambiguity. On the segment joining
any two reachable moment vectors, the denominator is a report probability
under a posterior mixture and is at least $\kappa$. Coefficients are
uniformly bounded and raw moments belong to $[0,1]$. Differentiating the
quotient gives a Lipschitz constant independent of all exponent gaps.

At every stage Theorem~\ref{thm:intrinsic-checkpoint} provides a raw
reachable codebook of cardinality at most $M$ and covering radius at most
$C_n\Psi_{n,N-n}(M,a)^{1/2}$. Update a representative by
\eqref{eq:formal-multiindex-update}, then quantize in the next reachable
codebook. Updating a reachable representative produces a reachable next
state, so this quantizer is always evaluated on its domain. If $e_n$
is the raw error, then
\[
 e_0=0,\qquad e_{n+1}\le L_ne_n+
        C_{n+1}\Psi_{n+1,N-n-1}(M,a)^{1/2}.
\]
The finite recurrence bounds every $e_n$ by
$C_N\Xi_N(M,a)^{1/2}$. The query matrix converts this to the
asserted squared-regret bound. The exact prefix appears only in this
error analysis, not in the transducer. No divided-difference coordinate,
reciprocal pivot or stored command tape is used in a transition.

Conversely, at each checkpoint the streaming index is an $M$-message
encoder of the prefix. All command laws used in
Theorem~\ref{thm:intrinsic-checkpoint} are prefixes of one common
exploration law. Each transducer therefore satisfies every checkpoint
lower bound; taking their maximum and then the infimum proves the lower
bound in \eqref{eq:intrinsic-streaming}. Finitely many checkpoint
constants can be replaced by their positive minimum. One checkpoint
and its remaining query trials are executed in a run, as before.
\end{proof}



\subsection{Confluent observation algebras and resolved directions}\label{sec:confluence}
The cardinality of a sumset is discontinuous when two of its elements
coalesce. The corresponding experiment need not be discontinuous. We now
identify the scales of the directions lost at such a collision. The
parameter $\theta$ below indexes a \emph{known calibration}; the unknown
parameter of each experiment is still $t$.

\subsubsection{Affine exponent families}
Fix $0=a_0<a_1<\cdots<a_{r-1}$ and real numbers $b_i$, with $b_0=0$.
On a sufficiently small compact interval $0\le\theta\le\theta_0\le1$ put
\[
 A_\theta=\{a_i+\theta b_i:0\le i<r\}.
\]
The order of these exponents is assumed unchanged. Take $I=[0,1]$ and
one fixed full-support prior $\mu$. A fixed coefficient matrix specifies
positive detector cells
\begin{equation}\label{eq:affine-calibration}
 k_j^\theta(t)=\sum_{i=0}^{r-1}c_{ji}t^{a_i+\theta b_i},
 \qquad \sum_jk_j^\theta=1.
\end{equation}
It has column rank $r$, and the cells have a common positive lower bound.
The comparator cube is fixed. These hypotheses include any sufficiently
small affine perturbation of a strictly positive calibration of full rank.
No continuity of an unknown prior density is assumed.

For a fixed positive future length $m\ge1$, let
\[
 \mathcal B_m=\left\{\left(\sum_{\nu=1}^m a_{i_\nu},
                              \sum_{\nu=1}^m b_{i_\nu}\right):
                              0\le i_\nu<r\right\}.
\]
Repeated \emph{pairs} are identified. Distinct pairs give distinct
exponents $\lambda+\theta\gamma$ for $0<\theta\le\theta_0$; shrink
$\theta_0$ to exclude crossings if necessary. For each first coordinate
$\lambda$, order its second coordinates as
$\gamma_{\lambda,0}<\cdots<\gamma_{\lambda,s_\lambda-1}$.
Set $K_m=\sum_\lambda s_\lambda$. The cluster at $\lambda=0$ consists
only of $(0,0)$.

Define the divided-difference functions
\begin{equation}\label{eq:divided-tests}
 \psi_{\lambda,k}^\theta(t)
  =[\lambda+\theta\gamma_{\lambda,0},\ldots,
       \lambda+\theta\gamma_{\lambda,k}]\,(x\mapsto t^x),
 \quad 0\le k<s_\lambda.
\end{equation}
At $\theta=0$ their continuous values are
\begin{equation}\label{eq:jet-tests}
 \psi_{\lambda,k}^0(t)=t^\lambda(\log t)^k/k!.
\end{equation}
For $\lambda>0$ these functions have value zero at $t=0$; the sole
zero-exponent function is the constant one. The bracket in
\eqref{eq:divided-tests} is a divided difference in the exponent, not a
new apparatus output. For a history $h$ let $z_\theta(h)$ be its posterior
expectations of these functions, omitting the constant coordinate.

List the orders $k$ of the nonconstant functions in nondecreasing order:
\[
 0=\nu_1\le\cdots\le\nu_d,\qquad d=K_m-1.
\]
There are $|mA_0|-1$ zero orders. The associated resolution function is
\begin{equation}\label{eq:resolution-function}
 \Phi_\theta(M)=\max_{1\le\ell\le d}
       \theta^{\,2(\nu_1+\cdots+\nu_\ell)/\ell}M^{-2/\ell},
       \qquad M\ge1.
\end{equation}
We set $0^0=1$ and $0^a=0$ for $a>0$ in this formula. In particular
$\Phi_0(M)=M^{-2/(|mA_0|-1)}$.

Choose $r$ attainable failure factors which form a basis in the formal
monomials $t^{a_i+\theta b_i}$, with coefficients independent of $\theta$.
For example use $1/2$ and $1/2+\delta t^{a_i+\theta b_i}$, $i>0$,
with $\delta$ sufficiently small. The fixed right inverse of the
calibration matrix realizes these by fixed command labels. All ordered
$m$-fold products form a common finite probe menu. Any other such fixed
basis, and its full product menu, gives equivalent bounds below.

\begin{theorem}[Uniform resolved checkpoint law]\label{thm:confluent-law}
Suppose
\begin{equation}\label{eq:full-future-condition}
 K_m-1\le n(r-1).
\end{equation}
Supply independent uniform commands for the first $n$ trials, then reveal
an independent uniform query from the preceding product menu. Let
$\overline R_{M;n,m}^\theta$ be the optimal unconditional squared-prediction
regret of an arbitrary $M$-message checkpoint encoder. Let
$R_{M;n,m}^{\theta,\max}$ be its minimax counterpart over admitted histories.
There are constants $0<c\le C<\infty$, independent of $M$ and $\theta$,
such that
\begin{equation}\label{eq:confluent-law}
 c\Phi_\theta(M)\le\overline R_{M;n,m}^\theta
       \le R_{M;n,m}^{\theta,\max}\le C\Phi_\theta(M),
 \qquad 0\le\theta\le\theta_0.
\end{equation}
The encoders may be Borel and randomized as in
Definition~\ref{def:finite-state}. At the common interior binomial command
tuple the normalized map $h\mapsto z_\theta(h)$ is a submersion of rank
$d$, including at $\theta=0$. Its local inverse bounds and an unconditional
minorization can be chosen uniformly in $\theta$.
\end{theorem}

The upper bound in this theorem is a checkpoint bound. It does not assume
that access to an exact prefix is free during repeated compression. Every
streaming filter is a checkpoint encoder, so the lower bound applies to
such filters as well. Section~\ref{sec:uniform} proves a matching genuinely
streaming bound across an entire singular family. Theorem~\ref{thm:attainable-checkpoint}
below removes the full-future restriction by truncating the resolution flag,
and Theorem~\ref{thm:general-uniform-streaming} proves the fixed-horizon
causal version for the affine class. When $m=0$ there is no nonconstant
future query and the prediction regret is zero; formula
\eqref{eq:resolution-function} is not used in that case.

\subsubsection{Confluent positivity and the physical prediction metric}
We give the details of the analytic and geometric steps in
Theorem~\ref{thm:confluent-law}.

The proof of Lemma~\ref{lem:confluent-positive} appears in the principal analytic inputs.

\begin{lemma}[Desingularization and observable scales]\label{lem:observable-scales}
The functions in \eqref{eq:divided-tests}, and the needed derivatives in
$\theta$, are uniformly bounded on $[0,1]\times[0,\theta_0]$.
The posterior vectors $z_\theta(h)$ therefore lie in a fixed box.
For the common product probe menu let $p_\theta(h)$ be its probability
vector. There is a fixed matrix $G$ of full column rank such that
\begin{equation}\label{eq:physical-metric}
 p_\theta(h)-p_\theta(\widetilde h)
    =G\operatorname{diag}(\theta^{\nu_i})
          (z_\theta(h)-z_\theta(\widetilde h)).
\end{equation}
In particular the squared norm on the left is bounded above and below
by positive constants times
$\sum_i\theta^{2\nu_i}|z_i(h)-z_i(\widetilde h)|^2$,
with constants independent of $\theta$.
\end{lemma}
\begin{proof}
The divided difference of order $k$ is the integral of
$f^{(k)}(\sum_jw_jx_j)$ over the standard $k$-simplex, whose volume is
$1/k!$. This identity follows by iterating the fundamental theorem of
calculus, starting with $(f(y)-f(x))/(y-x)=\int_0^1f'(x+s(y-x))\dd s$.
It also gives the repeated-node value. Here $f(x)=t^x$, so every
integrand is $t^x(\log t)^k$. At a positive cluster the exponent is
bounded below by some $a>0$ uniformly on the fixed parameter interval.
For every fixed $j$, $\sup_{0<t\le1}t^a|\log t|^j<\infty$.
The same domination applies after finitely many $\theta$ derivatives.
There are no higher divided differences at zero. Posterior expectations
of bounded functions are bounded independently of the history.

Newton's interpolation identity, proved by successive subtraction of
the interpolating polynomials, is exact at each node:
\begin{equation}\label{eq:newton-exponents}
 t^{\lambda+\theta\gamma_{\lambda,j}}
   =\sum_{k=0}^j\theta^k
      \left(\prod_{v<k}(\gamma_{\lambda,j}-\gamma_{\lambda,v})\right)
                  \psi_{\lambda,k}^\theta(t).
\end{equation}
The matrix multiplying $\operatorname{diag}(\theta^k)$ is fixed,
triangular within each cluster, and has nonzero diagonal. The standard
identities used here are classical divided-difference calculus
\cite{deBoor}; the integral proof makes the uniformity at coalescence
explicit.

The expansion matrix from the formal pair monomials in $\mathcal B_m$
to the probe products has full column rank. To see this, invert the fixed
one-step basis change and multiply it $m$ times. Every formal monomial
product is a linear combination of the probe products. Identifying equal
pairs is a surjective linear quotient, so the spanning statement survives
that identification. Both this expansion matrix and the triangular
matrix in \eqref{eq:newton-exponents} are independent of $\theta$.
Their product has full column rank. Subtracting two normalized posterior
vectors removes its constant column and gives \eqref{eq:physical-metric}.
The remaining columns are still independent. Their extreme singular
values give the two constants. The equality extends to $\theta=0$ by
continuity; vanishing weights then remove precisely the unresolved jets.
\end{proof}

\begin{lemma}[Uniform attainable patch with its evidence]\label{lem:uniform-patch}
Under \eqref{eq:full-future-condition}, there are $\rho,\beta>0$,
independent of $\theta$, such that the subprobability distribution of
$z_\theta$ on all-failure prefixes dominates $\beta$ times uniform
probability on a translate of $[-\rho,\rho]^d$.
\end{lemma}
\begin{proof}
Use the binomial factors $(1+c_it^{D_\theta})/2$ of
Lemma~\ref{lem:binomial-tangent}, with distinct small $c_i$ and
$D_\theta=a_{r-1}+\theta b_{r-1}$. Since the coefficient calibration is
fixed, one common interior command tuple implements this family for every
$\theta$. Its unnormalized tangent has $n(r-1)+1$ distinct monomials,
also at zero. At positive $\theta$, pairing with the $K_m$ distinct
future exponents has rank $K_m$ by Lemma~\ref{lem:mixed-moment}.
The divided-difference basis change is invertible there. At zero,
Lemma~\ref{lem:confluent-positive} gives the same rank directly, without
passing a rank assertion through a singular limit. The product itself
lies in the tangent. The normalization map
$v\mapsto v-p e_0v$ has kernel exactly $\spanop\{p\}$ on this full
image, as in \eqref{eq:normalized-derivative}. Thus the normalized jet
map has rank $d$ at every $\theta$, including zero.

Here are uniform local details. Let $q$ be the total number of command
coordinates and $g_*$ the common interior tuple. The matrices
$J_\theta=D_gz_\theta(g_*)$ are continuous in $\theta$ and surjective.
Their least row singular value has a positive minimum $\sigma$ on the
compact parameter interval. Choose an orthonormal basis $E_\theta$ of
the row space and a complementary orthonormal basis $F_\theta$ of its
kernel, and use coordinates $g=g_*+E_\theta u+F_\theta w$.
The derivative at zero of
\[
 (u,w)\longmapsto
       (z_\theta(g_*+E_\theta u+F_\theta w)-z_\theta(g_*),w)
\]
is block diagonal with invertible upper block, least singular value at
least $\sigma$, and lower block the identity. Its derivative norm and
its second derivative have common finite bounds on a fixed neighborhood.
This follows from Lemma~\ref{lem:observable-scales}, finite product
rules, and the prefix evidence lower bound $\kappa^n$.

Choose a ball inside the command cube on which the derivative differs
from its central value by at most half its least singular value. This
radius can be chosen uniformly. For a sufficiently small target ball,
the map obtained by subtracting the nonlinear remainder and applying the
inverse central derivative is a contraction of that command ball into
itself. The contraction theorem gives a differentiable inverse. Hence a
product of a $d$-cube and a $(q-d)$-cube of common positive side length
lies in its image, uniformly in $\theta$. All inverse Jacobian
determinants are bounded below by a positive constant: the forward
derivative is uniformly bounded, so its determinant is uniformly bounded
above. The orthogonal command-coordinate change has determinant of
absolute value one. No continuity of the chosen orthonormal frames in
$\theta$ is needed for these uniform bounds.

Uniform commands have density $(1-2\eta)^{-q}$. Their joint law with the
all-failure word has density
\[
 (1-2\eta)^{-q}\mu\!\left(\prod_{i=1}^nF_{g_i}\right)
       \ge (1-2\eta)^{-q}\eta^n.
\]
Change variables using the inverse just constructed and integrate over
the complementary cube. This gives the asserted minorization, with
$\beta>0$ the product of this density bound, the inverse-Jacobian bound,
and the two cube volumes. The convention for a zero-dimensional
complement is volume one. This is a subprobability bound retaining the
failure evidence, not conditioning away a rare prefix.
\end{proof}

\subsubsection{A scale-sensitive finite-state calculation}
\begin{lemma}[Rectangular quantization at every budget]\label{lem:rectangle}
For $a_1\ge\cdots\ge a_d>0$ set
\[
 e_M(a)=\max_{1\le\ell\le d}
                  \left(\frac{a_1\cdots a_\ell}{M}\right)^{1/\ell}.
\]
A box with these side lengths has an at-most-$M$-cell rectangular
partition of diameter at most $2\sqrt d\,e_M(a)$ in every cell.
The optimal mean squared distortion of its uniform distribution using
$M$ arbitrary centers is at least $c_d e_M(a)^2$, with $c_d>0$
depending only on $d$. Zero side lengths may be omitted.
\end{lemma}
\begin{proof}
Put $e=e_M(a)$ and take $N_i=\max\{1,\lfloor a_i/e\rfloor\}$
equal intervals in coordinate $i$. If $\ell$ coordinates satisfy
$a_i\ge e$, then
$\prod_iN_i\le\prod_{i\le\ell}a_i/e^\ell\le M$.
Every side length in a cell is at most $2e$: for $a_i/e\ge1$ use
$\lfloor a_i/e\rfloor\ge a_i/(2e)$, and otherwise use $a_i<e$.
Choose an actual reachable representative in each nonempty cell when
covering a reachable subset of the box. The diameter bound remains valid.

For the lower bound project onto the first $\ell$ coordinates. The
projected uniform density is $(a_1\cdots a_\ell)^{-1}$ on its box.
The union of $M$ balls of radius $r$ has probability at most
$M v_\ell r^\ell/(a_1\cdots a_\ell)$, where $v_\ell$ is the volume of
the unit ball. Taking
$r=(a_1\cdots a_\ell/(2Mv_\ell))^{1/\ell}$ shows that mean squared
distance is at least $r^2/2$. Projection cannot increase distance.
Apply this for every $\ell$ and take the smallest of the finitely many
positive dimensional constants. Centers need not lie in the box.
\end{proof}

\begin{proof}[Proof of Theorem~\ref{thm:confluent-law}]
Let $D_\theta=\operatorname{diag}(\theta^{\nu_i})$; this temporary
matrix notation is unrelated to the peak dimension $D_A(N)$.
By Lemma~\ref{lem:observable-scales}, $D_\theta z_\theta(h)$ lies in
a translated box of side lengths at most $2B\theta^{\nu_i}$ for a fixed
$B$. Lemma~\ref{lem:rectangle} partitions this box into at most $M$
cells with squared diameter at most $C\Phi_\theta(M)$. Choose a
reachable representative in each nonempty cell. At the checkpoint,
encode the cell of the exact history and decode its representative's
probe probabilities. These are probabilities in $[0,1]$. Equation
\eqref{eq:physical-metric} gives the uniform upper regret bound.
Finite cells with deterministic boundary conventions give Borel encoders.

For the lower bound first average private decoder coins conditional on
the message and query. Squared loss cannot increase, leaving at most
$M$ centers in probe space. Orthogonally project these centers onto the
affine space $p_*+\operatorname{range}(G)$; projection cannot increase
distance to a true prediction. Invert $G$ on its range. Its lower singular
bound reduces the problem to distance from at most $M$ arbitrary centers
in the scaled $z$ coordinates. The law there dominates $\beta$ times
the image of the cube from Lemma~\ref{lem:uniform-patch}; its side lengths
are $2\rho\theta^{\nu_i}$. The lower half of
Lemma~\ref{lem:rectangle} gives $c\Phi_\theta(M)$ after division by
the fixed number of queries. Encoder randomization cannot beat distance
to the closest center at a fixed history. Fresh public randomness can be
fixed and then averaged because it is independent of the entire prefix.
At zero, omit zero-weight coordinates in this argument. The same constants,
reduced if necessary over the finitely many possible dimensions, work
throughout the compact parameter interval. This proves both criteria and
the claimed submersion statement.
\end{proof}

Thus the relevant finite-resolution data are the multiplicities of additive
collisions and their confluent orders, not only the cardinality of the
nearby sumset. The result is a uniform comparison up to constants. It
neither asserts a new general quantization principle nor identifies a
universal leading distortion constant.


\subsection{Finite-state filtering after every report}\label{sec:streaming}
Fix $N\ge2$ and abbreviate $d_n=d_A(n,N-n)$ and $D_*=D_A(N)>0$.
The exact state from Theorem~\ref{thm:causal} will be denoted $s_n$.
A clock $n$, calibration, prior moments and a precomputed finite program
are read-only data. The charged resource is the persistent state between
successive input symbols.

The resource model and loss are defined in Section~\ref{sec:experiment}.

\begin{lemma}[Uniform finite-horizon transition bounds]\label{lem:lipschitz}
On the compact exact state sets there are constants $L_n<\infty$ such
that the update $T_n(s,g,x)$ satisfies
\[
 \norm{T_n(s,g,x)-T_n(\widetilde s,g,x)}
 \le L_n\norm{s-\widetilde s}
\]
for every admitted $g,x$ and $s,\widetilde s\in\cS_n$.
Every fixed finite probe prediction vector is likewise Lipschitz in the
state, with a finite constant independent of the history.
\end{lemma}
\begin{proof}
In factor coordinates an update before changeover appends the new factor
$f/\mu f$. The old factor coordinates are unchanged. All unnormalized
report probabilities lie in $[\kappa,1]$; consequently the normalized
factors lie pointwise in $[\kappa,1/\kappa]$. Their coordinates are bounded,
since they are continuous images of a compact cube and finite alphabet.
At changeover each needed moment is a ratio of polynomial expressions in
these coordinates and the fixed prior moments, with denominator the
integral of a product of normalized factors, bounded below by $\kappa^{n+1}$.
The same lower bound holds on line segments between factor tuples,
because each individual interpolated normalized factor is still at least
$\kappa$. Finite differentiation and the mean value theorem give a
uniform derivative bound there.

In moment coordinates the update is \eqref{eq:sparse-update}. Its
denominator is at least $\kappa$ at every state; on a line segment between
two such moment vectors the corresponding posterior mixture has the same
lower bound. Numerators, coefficients and moment coordinates range over
compact sets. The quotient rule therefore gives a uniform bound. Probe
predictions are linear in the relevant moment coordinates, or normalized
polynomial expressions in the factor coordinates with the same positive
denominator control. The argument gives their claimed Lipschitz constants.
All constants use only finite algebra and the supplied moments; they may
grow with $N$ or deteriorate with $\kappa$.
\end{proof}

\begin{theorem}[Finite-state streaming upper bound]\label{thm:stream-upper}
For every $M\ge1$ there is a deterministic $M$-state filter, updated after
every report, whose prediction regret at every checkpoint and every
admitted command-report history is at most
\begin{equation}\label{eq:stream-upper}
 C_N M^{-2/D_*},
\end{equation}
where $C_N<\infty$ depends on the fixed experiment, prior, probe lists
and total horizon, but not on $M$ or the input history. The representatives
and read-only transition functions are supplied existentially from the exact
model; this is not an effective synthesis theorem from finite-precision data.
\end{theorem}
\begin{proof}
Enclose $\cS_n\subset\R^{d_n}$ in $[-B_n,B_n]^{d_n}$. For $d_n>0$
partition each coordinate interval into
$k_n=\lfloor M^{1/d_n}\rfloor$ intervals. In each nonempty cell choose
one representative \emph{in} $\cS_n$. There are at most $k_n^{d_n}\le M$
representatives, and assigning a state to its cell representative has
error at most
\[
 \delta_n\le 2B_n\sqrt{d_n}/k_n
 \le4B_n\sqrt{d_n}\,M^{-1/d_n}
 \le C M^{-1/D_*}.
\]
Here $C$ is the maximum of the finite prefactors. Empty cells need no
representative. For $d_n=0$ use the unique state. Deterministic boundary
and tie conventions give Borel quantizers $Q_n$.

The filter stores only the index of a representative $\widehat s_n$.
Its transition is exactly
\begin{equation}\label{eq:transducer}
 \widehat s_{n+1}
   =Q_{n+1}\bigl(T_n(\widehat s_n,g_{n+1},x_{n+1})\bigr).
\end{equation}
Since each representative is reachable and every report has positive
probability, its exact image lies in $\cS_{n+1}$. Thus the quantizer is
always applied inside its domain and no infeasible moment update is used.
The finite representative choices and the functions in
\eqref{eq:transducer} are read-only program data. They may be obtained by
selecting points in nonempty cells of the compact state sets; no efficient global synthesis
or computability of unsupplied real prior moments is asserted.

Let $e_n=\norm{s_n-\widehat s_n}$. On the actual command-report stream,
Lemma~\ref{lem:lipschitz} gives
\[
 e_0=0,\qquad e_{n+1}\le L_ne_n+C M^{-1/D_*}.
\]
Set $b_0=0$ and $b_{n+1}=L_nb_n+C$. Induction proves
$e_n\le b_nM^{-1/D_*}$. This explicitly accumulates the errors from
\emph{all} previous lossy updates. At a query use the probe predictions
of $\widehat s_n$, which are actual probabilities. If the vector prediction
map has Lipschitz constant $K_n$, \eqref{eq:brier} bounds regret by
$(K_n^2/s_n^{\rm q})b_n^2 M^{-2/D_*}$. The maximum of these finitely
many constants proves \eqref{eq:stream-upper}. No exact prefix was kept
for the changeover: the factor-to-moment computation uses the representative
factors already encoded by $i_n$.
\end{proof}

\begin{theorem}[Matching exponent in one fixed streaming experiment]\label{thm:stream-lower}
Choose $n_*$ attaining $D_*=d_A(n_*,N-n_*)$. For the first $n_*$ trials
supply independent uniform commands from $\cG$, independent of the
parameter, cartridges and filtering coins. The realized command is read
once by the filter. At $n_*$ reveal the independent probe query and execute
its remaining trials. This is one fixed experiment, known before the
filter is chosen. Its optimal unconditional regret over $M$-state filters
satisfies
\begin{equation}\label{eq:stream-match}
 0<c_NM^{-2/D_*}\le\overline{\mathcal R}_{M,N}^{\rm stream}
                  \le C_NM^{-2/D_*}.
\end{equation}
The same power is sharp for the minimax streaming criterion taking the
maximum conditional regret over checkpoints and input histories. Fresh
private update and decoder randomization does not defeat the lower bound;
a fresh public coding seed independent of the complete prefix may also be
allowed. Thus for fixed $N$, the required persistent bits for this prediction
task satisfy
$B_*(\varepsilon)=\tfrac{D_*}{2}\log_2(1/\varepsilon)+O_N(1)$.
\end{theorem}
\begin{proof}
Write $n=n_*$, $d=D_*$, and $s=s_n^{\rm q}$. Theorem~\ref{thm:rank}
gives an interior command tuple $g_0$ where the probe map $p(g)$ on
all-failure prefixes has rank $d$. Select $d$ independent probe coordinates,
and let $E$ denote their coordinate projection, of operator norm one.
Choose complementary command coordinates $w$ so that
$g\mapsto(Ep(g),w(g))$ has an invertible derivative at $g_0$.
Let $\Phi$ be its local inverse. There is a compact product
$B_\rho^d\times W$ inside its domain, centered at the projected
prediction point, such that $\Phi$ stays in the command interior and
$|\det D\Phi|\ge j_0>0$. The complementary box $W$ has positive
volume, with unit-volume convention in zero dimensions.

If $k=Kn$ is the number of real command entries, their density is
$(1-2\eta)^{-k}$. The joint density of commands and an all-failure
prefix is this density times
$Z(g)=\mu\prod_iF_{g_i}\ge\eta^n$. Change variables with $\Phi$ and
integrate over $W$. The subprobability law of the projected predictions
on this prefix dominates $\beta$ times uniform probability on
$B_\rho^d$, with
\[
 \beta=\frac{\eta^n j_0|W|\,|B_\rho^d|}{(1-2\eta)^k}>0.
\]
In particular the lower bound retains the probability of the failure
prefix. It does not assign an atom to an exact real command or assume
a flat ball in the entire curved prediction image.

Every streaming filter is in particular an $M$-message encoder at this
checkpoint. Given any message, average decoder coins conditional on each
query; squared loss can only decrease. This gives at most $M$ prediction
vectors. Randomizing the encoder cannot beat the closest of these vectors
at a fixed history. Projecting their centers by $E$ cannot increase squared
distance. For a uniform point $X$ in a $d$-ball of radius $\rho$ and any
$M$ centers, the union-of-balls estimate gives
\[
 \Pr\{\min_i\norm{X-z_i}\le r\}\le
                    \min\{1,M(r/\rho)^d\}.
\]
Integrating its complementary tail yields
\[
 \E\min_i\norm{X-z_i}^2
 \ge\int_0^{\rho^2 M^{-2/d}}
          (1-Mt^{d/2}/\rho^d)\dd t
 =\frac d{d+2}\rho^2 M^{-2/d}.
\]
Multiply by $\beta/s$ to obtain $c_N>0$. Fresh public randomness can be
fixed and then averaged because it is independent of the prefix. Private
online coins may equivalently be drawn on an independent tape for this
proof; fixing them makes a deterministic transducer and the same bound
holds uniformly. That proof device supplies no tape to the implemented
filter. The prescribed exploration is independent of that tape, a condition
that is essential to this particular argument. An old command-generation
seed is not fresh coding randomness and remains charged.

Theorem~\ref{thm:stream-upper} gives the upper bound in the same experiment.
The minimax conditional criterion is at least its unconditional regret,
while the universal upper filter works at all checkpoints. Finally put
$M=2^B$ and solve the two inequalities. The Euclidean covering mechanism
is classical \cite{GrafLuschgy}; the identified sparse dimension and the
actual repeated-update construction are the experiment-specific inputs.
\end{proof}

\begin{theorem}[A separate finite-state control bound]\label{thm:control-bound}
For the same finite detector and horizon, suppose stage rewards
$r_n(g,x)\in[0,1]$ are continuous in $g$ for each report, and the objective
is their expected sum. There is an $M$-state feedback controller whose
loss relative to the full-history optimum is at most
\begin{equation}\label{eq:control-bound}
 C'_N M^{-1/D_*}.
\end{equation}
No matching lower power for arbitrary control tasks is asserted. An
additional finite payoff automaton must be counted as task state, multiplying
the number of states if it is included.
\end{theorem}
\begin{proof}
Exact dynamic programming on $\cS_n$ is valid by
Proposition~\ref{prop:tests}. With terminal value zero, its recursion is
\[
 V_n(s)=\max_{g\in\cG}\sum_{x\in\cX}p_x(s,g)
       \{r_{n+1}(g,x)+V_{n+1}(T_n(s,g,x))\}.
\]
Here $p_x$ is a linear moment functional or a normalized factor expression.
The positive denominator bounds, compact action set and continuous rewards
make the displayed objective continuous, and the maximum is attained.
By backward induction it is uniformly Lipschitz in $s$, since finite sums,
products of bounded Lipschitz functions, and the updates in
Lemma~\ref{lem:lipschitz} preserve that property. The maximum of functions
with a common Lipschitz constant is again Lipschitz. Denote by $H_n$ a
uniform Lipschitz constant of the action objective, before its maximization.

At each representative choose an exact maximizing action and store that
action in the finite output table. At the true state $s_n$, the loss from
using the action optimal at $\widehat s_n$ is at most $2H_ne_n$.
Update the representative by \eqref{eq:transducer}. The true exact state
is computed only for analysis along the controller's actual observations;
the implemented controller still stores only its index. Bellman telescoping
and conditional expectation give total loss at most
$2\sum_{n<N}H_ne_n\le2\sum_{n<N}H_nb_n M^{-1/D_*}$.
This proves \eqref{eq:control-bound}. It uses no continuity of the optimizing
action as a function of state, and it does not differentiate an optimizer.
Its approximation principle is related to \cite{Subramanian}; it is not
claimed as a new general dynamic-programming principle.
\end{proof}

\subsubsection{Horizon and calibration remain part of the constants}
If every factor in a probe menu satisfies $0<F\le v<1$, then
$H_{n,j}\le v^{N-n}$ and every posterior prediction satisfies the same
bound. The one-state decoder that always announces zero therefore has
\begin{equation}\label{eq:zero-memory}
 \mathcal R\le v^{2(N-n)}.
\end{equation}
For the shared physical menu $F=c+\delta k_j$ one may take $v=c+\delta$.
Thus a checkpoint implementation may choose, in advance, between the
streaming filter and the one-state decoder and obtain
\[
 \mathcal R\le\min\{C_N M^{-2/D_*},(c+\delta)^{2(N-n)}\}.
\]
At $c=1/2$, $\delta=1/8$, $N-n=20$, the second bound is
$(5/8)^{40}<10^{-8}$. This does not contradict a fixed-$N$ high-resolution
exponent. Neither the new sparse rank theorem nor the streaming construction
turns it into a horizon-uniform useful-memory law. We have not rescaled
the loss or discarded rare histories to avoid this bound.

The upper constants involve state bounds, transition Lipschitz constants
and repeated error propagation. The lower constants involve an attainable
Jacobian, a chart radius, evidence and the prior pairing. They are not
uniform over vanishing detector amplitudes, all full-support priors or
increasing horizons. With a finite fixed command alphabet and finite $N$
there are finitely many histories and an exact finite code eventually
suffices; the matching asymptotic lower bound above expressly uses the
fixed continuous-command exploration protocol.


\subsection{Attainable resolution flags and uniform causal memory}\label{sec:attainable-flag}
The full-future hypothesis in Section~\ref{sec:confluence} makes every
nonconstant future jet attainable. It is not needed for a resolution law,
but it cannot simply be omitted from the formula there. When the past is
shorter, the relevant list is an initial part of the future resolution
flag. The lower bound requires transversality to this flag; the upper
bound requires an estimate on the whole, generally curved, attainable
image. We establish both and then prove their compatibility with every
successive report update.

\subsubsection{Statement of the attainable law}
Retain the affine calibration and positivity assumptions of
Section~\ref{sec:confluence}, but do not impose
\eqref{eq:full-future-condition}. Fix $n,m\ge1$. There are $K_m$ formal
sum pairs in $\mathcal B_m$; the zero pair is the constant, and the other
$q_m=K_m-1$ divided-difference coordinates have orders
\[
 \nu_{m,1}\le\cdots\le\nu_{m,q_m}.
\]
Order equal orders by increasing limiting exponent, once and for all.
Thus every initial part of the list includes, within each cluster, all
jets preceding any jet that it contains. Set
\begin{equation}\label{eq:attainable-orders}
 p_{n,m}=\min\{n(r-1),q_m\},\qquad
 S_{m,\ell}=\sum_{i=1}^{\ell}\nu_{m,i},\qquad
 \Phi^{\rm att}_{n,m}(M,\theta)
 =\max_{1\le\ell\le p_{n,m}}
       \theta^{2S_{m,\ell}/\ell}M^{-2/\ell}.
\end{equation}
Here $0^0=1$ and $0^s=0$ for $s>0$. If $n=0$ or $m=0$, define
$\Phi^{\rm att}_{n,m}=0$ separately. No division by a zero dimension is
then involved. Throughout, constants may depend on the fixed experiment,
prior, horizon and calibration interval, but not on $M$ or $\theta$.

\begin{theorem}[Attainable confluent checkpoint law]\label{thm:attainable-checkpoint}
For every integer $M\ge1$ and every calibration in the compact interval
of Section~\ref{sec:confluence}, the unconditional optimal $M$-message
checkpoint regret under independent uniform commands and the minimax
checkpoint regret both have order
$\Phi^{\rm att}_{n,m}(M,\theta)$, with positive uniform comparison
constants. The upper code has at most $M$ reachable representatives
and works on every admitted history. The lower bound permits the
independent encoder and decoder randomization specified in
Definition~\ref{def:finite-state}.

In particular, past limitation truncates the ordered ambient list at
$p_{n,m}$; it does not introduce the remaining ambient orders into the
rate. At zero the exponent is
$2/\min\{n(r-1),|mA_0|-1\}$, and at each fixed positive calibration it
is $2/p_{n,m}$.
\end{theorem}

For the sequential assertion fix $N\ge2$ and choose the compact
calibration interval so that distinct formal pairs in
$\mathcal B_m$, $1\le m<N$, do not meet at a positive calibration.
Such an interval always exists after decreasing its positive endpoint:
there are finitely many pairs and their exponent differences are affine.
Put
\begin{equation}\label{eq:general-stream-profile}
 \Omega_N(M,\theta)
   =\max_{1\le n<N}\Phi^{\rm att}_{n,N-n}(M,\theta).
\end{equation}
Use the same independent uniform command law at every acquisition time.
For each selected checkpoint reveal its independent product query only
after the index has been formed. As in Section~\ref{sec:streaming}, only
one checkpoint and its remaining query trials are executed in a run.
Let $\mathfrak R^{\rm av}_{M,N}$ be the infimum, over one streaming
filter, of the maximum of its unconditional checkpoint regrets. Let
$\mathfrak R^{\rm max}_{M,N}$ be the corresponding infimum of the
maximum conditional regret over checkpoints and histories.

\begin{theorem}[Uniform causal realization of the resolution flag]\label{thm:general-uniform-streaming}
In this fixed-horizon experiment there are $0<c\le C<\infty$ such that
\begin{equation}\label{eq:general-stream-law}
 c\Omega_N(M,\theta)\le\mathfrak R^{\rm av}_{M,N}(\theta)
 \le\mathfrak R^{\rm max}_{M,N}(\theta)
 \le C\Omega_N(M,\theta)
\end{equation}
for every $M\ge1$ and every admitted $\theta$, including zero. For each known $\theta$ and each $M$, there exists a
single deterministic $M$-state transducer attaining the upper bound after
every report. Its codebooks and transitions may depend on $\theta$ and
$M$; the comparison constants do not. No joint-in-calibration measurable
or computable selection of codebooks is asserted. It stores only a representative index; its transitions use
physical moments, never division by a collision parameter. Read-only
prior moments at the formal exponents in $\mathcal B_N$ suffice.
The lower bound holds in the stated common exploration experiment,
including for the independent randomized filters of
Definition~\ref{def:finite-state}.
\end{theorem}

These are statements about finite persistent memory in the exact-real
program model already defined, not about finite-precision synthesis.
In particular the number of real coordinates used to describe a
representative is not the persistent state count: the latter is the
number of representative labels, at most $M$.

\subsubsection{A global covering estimate}
We first record the real-geometric input with its anisotropic consequence.
A bounded semialgebraic format means a bound on the number of variables,
polynomial relations and their degrees in a defining formula; arbitrary
real coefficients are allowed. Quantified formulas of fixed format can
also be used, since real quantifier elimination gives a bound depending
only on this format. Neither a coefficient norm nor a lower bound on a
nonzero coefficient is part of the definition.

The proof of Lemma~\ref{lem:tame-rectangle} appears in the principal analytic inputs.

The bounded-format hypothesis is essential to a uniform statement of
this kind. An unrestricted smooth image could develop arbitrarily many
sheets. In the present experiment this hypothesis is a consequence of
the finite product likelihood, not an extra regularity assumption on the
prior or on an unknown attainable variety.

\begin{lemma}[The whole attainable image has bounded format]\label{lem:attainable-format}
Let $z_\theta(h)\in\R^{q_m}$ be the nonconstant divided-difference
moment vector, in the ordered coordinates of
\eqref{eq:attainable-orders}, and put
\[
 D_{m,\theta}=\operatorname{diag}(\theta^{\nu_{m,i}}),\qquad
 S_{n,m}^\theta=\{D_{m,\theta}z_\theta(h):|h|=n\}.
\]
The set $S_{n,m}^\theta$ is compact, has semialgebraic format bounded
independently of $\theta$, and has dimension at most $p_{n,m}$.
It is contained in a rectangle with sides
$[-B\theta^{\nu_{m,i}},B\theta^{\nu_{m,i}}]$, with $B$ uniform.
At zero all positive-order coordinates are zero.
\end{lemma}
\begin{proof}
Fix the finite report word $x_1,\ldots,x_n$ and write $u$ for its
command entries. Every report factor is affine in its current command.
Thus the evidence
$Z(u)=\mu\prod_{i=1}^n f_{x_i,g_i}^\theta$ and each numerator
\[
 N_j(u)=\mu\left(\psi_j^\theta
                   \prod_{i=1}^n f_{x_i,g_i}^\theta\right)
\]
are polynomials of degree at most $n$ in $u$. Their coefficients are
real prior integrals, well defined also at zero. The moment image is
specified by the existence of $u$ in its fixed cube with
$Z(u)y_j=\theta^{\nu_{m,j}}N_j(u)$ for every $j$.
Moreover $Z(u)\ge\kappa^n>0$. This is a semialgebraic formula of
fixed format, uniformly in its real coefficients. Finite union over the
report words preserves a format bound, and continuity on compact command
cubes proves compactness. There is no requirement that these real
coefficients, as functions of $\theta$, be semialgebraic.

For the dimension bound, divide each report factor by its prior
integral. Its coefficient vector lies in the affine hyperplane
$\mu f=1$ of the $r$-dimensional one-step space. The constant monomial
allows its constant coefficient to be eliminated, leaving at most $r-1$
coordinates. The product posterior and all displayed moments are rational
functions of these $n(r-1)$ coordinates. A rational map cannot increase
semialgebraic dimension. Multiplication by $D_{m,\theta}$ cannot increase
it either; the ambient dimension is $q_m$. Finally the uniform bounds
for $\psi_j^\theta$ in Lemma~\ref{lem:observable-scales} give
$|z_j|\le B$ before any command integration: they are expectations under
probability measures. This proves the rectangle assertion without an
inverse power of the prefix evidence.
\end{proof}

\subsubsection{Transversality to every initial resolution flag}
\begin{lemma}[Attainment of the truncated flag]\label{lem:attained-flag}
At one common interior binomial command tuple, projection of
$z_\theta(h)$ onto its first $p_{n,m}$ coordinates has full row rank
for every calibration, including zero. Under independent uniform commands,
the all-failure subprobability distribution of this projection dominates
$\beta$ times uniform probability on a translated cube of side $2\rho$,
where $\beta,\rho>0$ are independent of the calibration.
\end{lemma}
\begin{proof}
Choose distinct sufficiently small $c_1,\ldots,c_n$ and realize the
factors $(1+c_it^{D_\theta})/2$ by the fixed right inverse of the
calibration matrix, as in Lemma~\ref{lem:interior}. Their commands stay
in a common interior neighborhood. Lemma~\ref{lem:binomial-tangent}
identifies their unnormalized tangent with a monomial space of dimension
$n(r-1)+1$, for every calibration, including zero.

Adjoin the constant to the selected $p=p_{n,m}$ future tests. Within
each limiting exponent cluster, the selected tests are exactly an initial
block of its divided differences. At a positive calibration this block
is an invertible linear transformation of the monomials at its first
nodes. Their union consists of distinct monomials. The mixed pairing
with the binomial tangent therefore has column rank $p+1$, by
Lemma~\ref{lem:mixed-moment}. At zero these same blocks become complete
initial logarithmic-jet blocks. Lemma~\ref{lem:confluent-positive}
gives column rank $p+1$ directly at zero. This is not an inference of
rank at zero from rank on a punctured interval.

Let $L$ be this pairing, let $F$ be the product at the witness, and
write $v=LF/\mu F$, with constant coordinate $e_0v=1$.
The normalized differential has image
\[
 (\mu F)^{-1}\{LQ-v\,e_0LQ:Q\hbox{ in the product tangent}\}.
\]
Since $L$ is onto $\R^{p+1}$, the bracket is precisely the
$p$-dimensional hyperplane $e_0=0$. Deleting the constant proves the
asserted row rank. The remaining, unselected future coordinates need not
be independent and are not included in this submersion assertion.

For clarity we give the uniformity and probability argument. On the
compact calibration interval, the least row singular value of the
selected derivative has a positive minimum, by the rank just proved
and continuity of all divided-difference derivatives. The second command
derivatives have a common bound, since the factors stay positive and
the test functions are uniformly bounded. Complete the row coordinates
with orthonormal kernel coordinates. On a fixed small command ball the
resulting square derivative differs from its value at the center by at
most half its uniformly positive least singular value. The contraction
proof of the inverse theorem gives a common product box in the selected
output and complementary coordinates. Shrink it so its inverse stays
in the command interior. Uniform forward derivative bounds give a
positive lower bound for the absolute inverse Jacobian. These estimates
use only orthogonality of the chosen frames; a continuous choice of
frames is unnecessary.

If $k$ is the number of command entries, their joint density is
$(1-2\eta)^{-k}$. The density of commands together with the all-failure
word is this density multiplied by its evidence, at least $\eta^n$.
Change variables using the local inverse and integrate over the common
complementary box. The resulting measure dominates a fixed positive
multiple of Lebesgue measure on the selected $p$-cube. Absorb the cube
volume into $\beta$. This proves the stated subprobability minorization;
it does not condition away the probability of the failure word. A
full-support prior without a density is permitted throughout, since the
continuous density used here is that of the commands.
\end{proof}

\begin{proof}[Proof of Theorem~\ref{thm:attainable-checkpoint}]
The physical prediction identity of Lemma~\ref{lem:observable-scales}
reads, after the fixed coordinate permutation,
\begin{equation}\label{eq:flag-query-metric}
 p_\theta(h)-p_\theta(\widetilde h)
 =G_mD_{m,\theta}\bigl(z_\theta(h)-z_\theta(\widetilde h)\bigr),
\end{equation}
where $G_m$ has full column rank and is independent of $\theta$.
Apply Lemmas~\ref{lem:tame-rectangle} and
\ref{lem:attainable-format} to the complete attainable image. Since
$0\le\theta\le1$, its ordered side lengths are
$B\theta^{\nu_{m,i}}$. Equation~\eqref{eq:tame-cover-radius} supplies
at most $M$ reachable representatives with squared covering radius
at most $C\Phi^{\rm att}_{n,m}$. The fixed norm of $G_m$ and the
finite query average convert this into the required uniform prediction
upper bound. This is a cover of all report strata, not just the local
patch used in the lower bound.

For that lower bound, first average decoder randomization for each
message and query. The outputs form at most $M$ centers. Orthogonally
project them to the affine span $p_*+\operatorname{im}G_m$ and apply
the fixed left inverse of $G_m$. Distances can decrease only by a
fixed factor. Project further onto the first $p_{n,m}$ coordinates.
Lemma~\ref{lem:attained-flag} then gives a subprobability multiple of
uniform measure on a translated rectangle of side lengths
$2\rho\theta^{\nu_{m,i}}$, $i\le p_{n,m}$. The lower half of
Lemma~\ref{lem:rectangle}, projecting onto each first $\ell$ axes,
gives $c\Phi^{\rm att}_{n,m}$. Centers may be outside the image;
randomized assignment cannot improve on the nearest center. Fixing and
averaging an independent public coding seed gives the same bound.
At zero, omit the zero-width axes. There are
$\min\{n(r-1),|mA_0|-1\}$ surviving axes, which also proves the
stated limiting exponent. For positive calibration the last power
$M^{-2/p_{n,m}}$ dominates as $M\to\infty$. The unconditional lower
bound is also a minimax lower bound, completing the proof.
\end{proof}

\subsubsection{Physical moments and successive updates}
For $b=(\lambda,\gamma)\in\mathcal B_m$, write
\[
 v_b(h)=\mu(P_h t^{\lambda+\theta\gamma}),\qquad v_{(0,0)}=1.
\]
We keep the formal labels even if the functions coincide at zero.
Let $v^{(m)}$ denote the nonconstant vector. Newton interpolation within
each cluster, as in Lemma~\ref{lem:observable-scales}, gives a fixed
invertible matrix $T_m$ with
\begin{equation}\label{eq:raw-versus-scaled}
 v^{(m)}=T_mD_{m,\theta}z_\theta.
\end{equation}
The constant cluster has no other member, so no deleted constant mixes
with this equation. Both $T_m$ and its inverse have fixed finite norms.
At zero the equation still holds with the appropriate repeated raw
moments. It is the analysis of this equation, not inversion of its
vanishing diagonal, that identifies the metric scales.

\begin{lemma}[Uniform updates in the physical metric]\label{lem:raw-uniform-update}
The map from reachable raw moments at remaining length $m$ to those at
length $m-1$, after a command and a report, is uniformly Lipschitz in
the raw Euclidean metric, independently of calibration. Every reachable
representative maps to a reachable next vector.
\end{lemma}
\begin{proof}
Write $e_i=(a_i,b_i)$ and the current report factor as
$f=\sum_i f_i(g,x)t^{a_i+\theta b_i}$. Its coefficients are uniformly
bounded, and $f\ge\kappa>0$. Addition of formal pairs gives
$b+e_i\in\mathcal B_m$ whenever $b\in\mathcal B_{m-1}$.
Consequently the exact update is
\begin{equation}\label{eq:raw-causal-update}
 v'_b=\frac{\sum_i f_i(g,x)v_{b+e_i}}
                  {\sum_i f_i(g,x)v_{e_i}},
       \qquad b\in\mathcal B_{m-1}.
\end{equation}
Every denominator is the posterior expectation of $f$, hence at least
$\kappa$. A line segment between reachable raw vectors represents a
mixture of their posterior probability measures, so the same bound holds
along the segment. All numerators, coefficients and coordinates have
uniform bounds. The quotient rule and the mean value theorem prove the
Lipschitz assertion. No exponent separation occurs in a denominator.
The formula is valid also when formal labels have equal values. Finally
a history realizing the representative, followed by the specified input,
realizes its updated vector; positivity admits every report.
\end{proof}

\begin{proof}[Proof of Theorem~\ref{thm:general-uniform-streaming}]
For each $1\le n<N$, Lemmas~\ref{lem:tame-rectangle} and
\ref{lem:attainable-format}, followed by the fixed matrix $T_{N-n}$,
give a raw-moment quantizer $Q_n$ with at most $M$ reachable
representatives and radius
\[
 \delta_n\le C_n\sqrt{\Phi^{\rm att}_{n,N-n}(M,\theta)}
             \le C_0\sqrt{\Omega_N(M,\theta)}.
\]
The maximum $C_0$ is over finitely many stages. At time zero the prior
moment vector is one fixed state; at time $N$ the remaining state is
constant. Initialize at the prior vector and implement
\[
 \widehat v_{n+1}=Q_{n+1}
       \bigl(U_n(\widehat v_n,g_{n+1},x_{n+1})\bigr),
\]
where $U_n$ is \eqref{eq:raw-causal-update}. The transition reads the
old representative index and the current input only. Reachability makes
its argument an element of the next quantizer's domain. The finite
codebooks and these real-valued functions are read-only program data.

Let $e_n=\|v_n-\widehat v_n\|$. With $e_0=0$,
Lemma~\ref{lem:raw-uniform-update} yields
\[
 e_{n+1}\le L_ne_n+C_0\sqrt{\Omega_N(M,\theta)}.
\]
Set $b_0=0$ and $b_{n+1}=L_nb_n+C_0$. Induction gives
$e_n\le b_n\sqrt{\Omega_N(M,\theta)}$ after every update.
Each query prediction is a linear form with fixed coefficients in the
raw formal moments. Its Lipschitz constant is therefore uniform, and
\eqref{eq:brier} proves the last inequality in
\eqref{eq:general-stream-law}. There is no exact-prefix recomputation:
all intermediate vectors in this implementation are functions of a finite
representative label and the current command-report pair.

For every streaming filter, its index at checkpoint $n$ is an
$M$-message encoder of the prefix. Theorem~\ref{thm:attainable-checkpoint}
therefore bounds that checkpoint's unconditional regret below by
$c_n\Phi^{\rm att}_{n,N-n}$. The commands in all these assertions
are prefixes of the same independent uniform acquisition rule. Taking
the maximum over $n$, and then the infimum over filters, gives the
first inequality with $c=\min_n c_n>0$. Independent filtering coins
can be fixed for this argument because they do not generate the
acquisition commands; averaging restores the randomized statement.
Finally a maximum conditional risk dominates its unconditional average,
which proves the middle inequality.
\end{proof}

\begin{corollary}[A uniform bit profile]\label{cor:general-bit-profile}
Let $B_N(\varepsilon,\theta)$ be the least number of persistent bits
needed to make the minimax streaming regret at most $\varepsilon$.
For $0<\varepsilon<1$, uniformly in the admitted calibration,
\begin{equation}\label{eq:general-bit-profile}
 B_N(\varepsilon,\theta)=
 \max\left\{0,\max_{\substack{1\le n<N\\1\le\ell\le p_{n,N-n}}}
 \left(\frac\ell2\log_2\frac1\varepsilon
       -S_{N-n,\ell}\log_2\frac1\theta\right)\right\}+O(1).
\end{equation}
At zero retain only terms with $S_{N-n,\ell}=0$. The same formula
holds for the maximum unconditional criterion.
\end{corollary}
\begin{proof}
For positive $\theta$, the condition
$\Omega_N(M,\theta)\le\varepsilon$ is equivalent to all inequalities
$M\ge\theta^{S_{N-n,\ell}}\varepsilon^{-\ell/2}$.
Taking the largest right-hand side, adjoining $M\ge1$, and then taking
base-two logarithms gives the expression. The two fixed comparison
constants in \eqref{eq:general-stream-law} change it by a bounded
amount, since the possible $\ell$ are bounded by the fixed horizon.
Rounding to a power of two changes it by at most one further bit.
The zero-calibration statement follows by omitting vanishing terms.
\end{proof}

\subsubsection{A past-limited peak with two weak directions}
The next consequence answers the concrete seven-trial boundary in the
source-pinned v7 technical note \cite{A1v7note}. It is a consequence
of the attainable-image argument, rather than an extrapolation of the
full-future formula.

\begin{corollary}[Seven-trial resolved memory]\label{cor:seven-trial}
For the cells \eqref{eq:near-collision-cells}, their fixed command cube
and the product probes \eqref{eq:common-near-probes}, take total horizon
seven. For the two streaming criteria above, uniformly for
$0\le\theta\le1/2$ and $M\ge1$,
\begin{equation}\label{eq:seven-law}
 \mathfrak R_{M,7}(\theta)\asymp
    \max\{M^{-1/3},\theta^{1/2}M^{-1/4}\}.
\end{equation}
The lower bound already holds at checkpoint four under the same uniform
exploration. The crossover budget has order $\theta^{-6}$ and the
crossover regret has order $\theta^2$.
\end{corollary}
\begin{proof}
Put $D=2+\theta$. For positive $\theta\le1/2$, the ten distinct
three-fold exponents are
\[
 0,\ 1,\ 2,\ D,\ 3,\ D+1,\ D+2,\ 2D,\ 2D+1,\ 3D.
\]
At zero their limiting clusters have multiplicities
$1,1,2,2,2,1,1$ at exponents $0,1,2,3,4,5,6$.
Thus the nonconstant order list has six zeros and three ones.
At $n=4,m=3$, however, $p_{4,3}=8$, and the selected list has six
zeros and only two ones. Theorem~\ref{thm:attainable-checkpoint} gives
\[
 \Phi^{\rm att}_{4,3}
 =\max\{M^{-1/3},\theta^{2/7}M^{-2/7},
                         \theta^{1/2}M^{-1/4}\}
 =\max\{M^{-1/3},\theta^{1/2}M^{-1/4}\}.
\]
The middle term is the product of the first term to power $3/7$ and
the last to power $4/7$, so cannot exceed their maximum. The minorized
eight-dimensional flag uses, at zero, the tests
\[
 1,t,t^2,t^2\log t,t^3,t^3\log t,t^4,t^5,t^6
\]
paired with the binomial tangent $\mathcal P_8$. Strict confluent
positivity proves its rank for every full-support prior. The third
ambient weak jet is not counted as an additional attainable dimension.

For completeness the extension of the upper bound to the whole stated
calibration interval does not require excluding positive collisions in
longer, unused sumsets. At stages $n=1,2,3$, the raw reachable images
have dimension at most $2n\le6$ and bounded format, and all raw
coordinates lie in $[0,1]$. Lemma~\ref{lem:tame-rectangle} gives
squared covering errors at most $CM^{-1/3}$ uniformly, whether or not
some formal raw coordinates coincide. At stage four use the preceding
anisotropic cover, valid throughout $[0,1/2]$ because its ten exponents
are distinct at every positive calibration in that interval. At stage
five, $m=2$ and the full-future cover has squared radius at most
$C\max\{M^{-1/2},\theta^{2/5}M^{-2/5}\}\le CM^{-1/3}$.
Stage six has two coordinates and squared covering radius at most
$CM^{-1}$. Formula~\eqref{eq:raw-causal-update} remains valid with
all formal labels at every stage, even when values coincide. Its uniform
Lipschitz estimate and the repeated-quantization argument prove the upper
bound in \eqref{eq:seven-law} on the entire interval. The stage-four
checkpoint lower bound proves the other direction.

Equating the two terms in \eqref{eq:seven-law} gives
$M=\theta^{-6}$ and common order $\theta^2$. These are comparison
orders, not a numerical integer threshold. At fixed positive calibration
the rate is $M^{-1/4}$, consistently with eight attainable dimensions;
the ambient nine-dimensional rate $M^{-2/9}$ does not occur.
\end{proof}
